\documentclass[12pt,onecolumn,letterpaper]{article}

\usepackage{geometry}
\geometry{letterpaper, margin=1in}

\usepackage{fontspec}
\usepackage{microtype}
\usepackage{setspace}
\usepackage{array}
\usepackage{booktabs}
\usepackage{longtable}
\usepackage{tabularx}
\usepackage{hyperref}
\usepackage{amsmath}
\usepackage[table]{xcolor}
\usepackage{caption}
\usepackage{enumitem}
\usepackage{etoolbox}
\usepackage{verbatim}
\usepackage{titlesec}
\usepackage{fancyhdr}
\usepackage{indentfirst}

\IfFontExistsTF{Cambria}{\setmainfont{Cambria}}{\setmainfont{TeX Gyre Termes}}
\IfFontExistsTF{Liberation Mono}{\setmonofont{Liberation Mono}[Scale=MatchLowercase]}{\setmonofont{Courier New}[Scale=MatchLowercase]}

\setlength{\parindent}{0pt}
\setlength{\parskip}{0.45em}
\setlength{\headheight}{15pt}
\setlength{\emergencystretch}{2em}
\setstretch{1.08}
\widowpenalty=10000
\clubpenalty=10000
\displaywidowpenalty=10000

\definecolor{RSBlue}{HTML}{16324F}
\definecolor{RSGray}{HTML}{E8EDF2}
\definecolor{RSRule}{HTML}{D6DCE3}

\titleformat{\section}{\Large\bfseries}{\thesection}{0.65em}{}
\titleformat{\subsection}{\large\bfseries}{\thesubsection}{0.6em}{}
\titleformat{\subsubsection}{\normalsize\bfseries}{\thesubsubsection}{0.5em}{}
\titlespacing*{\section}{0pt}{1.5em}{0.7em}
\titlespacing*{\subsection}{0pt}{1.1em}{0.45em}
\titlespacing*{\subsubsection}{0pt}{0.9em}{0.35em}
\renewcommand{\refname}{Referințe}

\pagestyle{plain}

\urlstyle{same}
\hypersetup{
    colorlinks=false,
    pdftitle={RomânăSimplă: propunere de normalizare și formalizare a limbii române},
    pdfauthor={Dan Micșa, PhD},
    pdfsubject={Propunere de normalizare și formalizare pentru nucleul RomânăSimplă}
}

\captionsetup[table]{font=small,labelfont=bf,labelsep=period,skip=6pt,justification=raggedright,singlelinecheck=false}

\newcommand{\rs}[1]{\texttt{#1}}
\newcolumntype{Y}{>{\raggedright\arraybackslash}X}
\newcommand{\RStablelayout}{%
\footnotesize
\setlength{\tabcolsep}{6pt}%
\renewcommand{\arraystretch}{1.12}%
\arrayrulecolor{RSRule}%
}
\newcommand{\RSDenseTablelayout}{%
\footnotesize
\setlength{\tabcolsep}{4pt}%
\renewcommand{\arraystretch}{1.03}%
\setlength{\LTleft}{\fill}%
\setlength{\LTright}{\fill}%
\arrayrulecolor{RSRule}%
}
\newcommand{\RSEBNFLayout}{\footnotesize}

\title{\Large\textbf{RomânăSimplă: O propunere de normalizare și formalizare a limbii române}}

\author{\textbf{Dan Micșa, PhD}\\
\small \href{mailto:dmicsa@gmail.com}{dmicsa@gmail.com}\\
\small \url{https://github.com/dmicsa/Rom-n-Simpl-.git}}
\date{Martie 2026}

\begin{document}

\maketitle
\thispagestyle{empty}

\begin{quote}
\small\textit{AI Attribution: Acest document a fost elaborat cu asistență extinsă din partea modelelor lingvistice mari. \textbf{GPT-5.4} a contribuit la redactare, reorganizare academică, formalizare și integrarea exemplelor tehnice în forma publicată aici. Responsabilitatea intelectuală pentru direcția proiectului, pentru deciziile de normalizare și pentru validarea conținutului rămâne a autorului uman.}
\end{quote}

\clearpage

\begin{abstract}
RomânăSimplă este o propunere de normalizare și formalizare controlată a limbii române,
construită în jurul a patru criterii explicite: \textit{Compact},
\textit{Precis}, \textit{Simplu} și \textit{Uniform}. Nucleul actual urmărește
o gramatică mică și testabilă, în care timpul este scris înaintea verbului,
forma verbală rămâne stabilă, definitudinea se marchează numai când este
necesar prin \rs{@}, posesia prin \rs{\#}, iar propozițiile grele sunt
descompuse în unități scurte. Lucrarea prezintă proiectul ca răspuns la
complexitatea gramaticală a românei standard și la densitatea ei de excepții,
propunând un nucleu disciplinat, explicit și formalizabil. Documentul sintetizează
specificația, corpusul canonic și formalizarea gramaticală într-o formă orientată
spre publicare. Rezultatul este o propunere coerentă de normalizare pentru un
subansamblu controlat al limbii, însoțită de criterii de acceptare, exemple
validate, operatori descriși sistematic și anexe tehnice. În forma actuală,
proiectul rămâne deliberat limitat la un nucleu restrâns, nu la o gramatică
exhaustivă, tocmai pentru a păstra testabilitatea și disciplina formală.
\end{abstract}

\tableofcontents
\vspace{1em}

\section{Introducere}

RomânăSimplă pornește de la o ipoteză simplă: o parte importantă din costul
românei standard stă nu în vocabularul de bază, ci în flexiune, în marcaje
redundante și în variații care mută informația în prea multe locuri simultan.
Proiectul nu urmărește o reformă totală a limbii, ci o variantă controlată,
comparabilă și testabilă, potrivită pentru transcriere simplificată, analiză
formală și evaluare pe corpus.

Forma actuală a sistemului este intenționat restrânsă. Nucleul nu încearcă să
acopere toate funcțiile românei contemporane, ci fixează un set mic de reguli
care pot fi aplicate repetabil. Accentul cade pe stabilitate și pe verificare,
nu pe expresivitate maximă.

\section{Cadru metodologic}

\subsection{Filtrul CPSU}

Toate regulile active sunt judecate prin filtrul CPSU: \textit{Compact},
\textit{Precis}, \textit{Simplu}, \textit{Uniform}. O regulă nouă intră în
nucleu numai dacă îmbunătățește clar cel puțin două dintre aceste criterii și nu
deteriorează sever celelalte două.

\begin{table}[h!]
\centering
\RStablelayout
\begin{tabularx}{\linewidth}{@{}lYY@{}}
\rowcolor{RSGray}
\textbf{Principiu} & \textbf{Întrebare de control} & \textbf{Consecință operațională} \\
Compact & Spune aceeași informație în mai puține semne? & Se preferă forme scurte, analitice și propoziții cu un singur fenomen dominant. \\
Precis & Evită deducția riscantă? & Nu se inventează timp, număr sau definitudine absente din sursă. \\
Simplu & Scade costul de învățare? & Se limitează numărul de reguli active și de excepții. \\
Uniform & Se aplică regula la fel în tot corpusul? & Se fixează lexic, ordine și formate stabile. \\
\bottomrule
\end{tabularx}
\caption{Filtrul CPSU folosit pentru selecția regulilor active}
\end{table}

Două reguli auxiliare traversează toate deciziile proiectului. Prima este
preferința pentru forma nemarcată: dacă două formulări păstrează aceeași
informație de bază, nucleul alege varianta cu mai puține marcaje vizibile. A
doua este regula \textit{operator minim necesar}: operatorii nu se introduc doar
pentru simetrie grafică, ci numai când aduc un contrast real.

\subsection{Surse și material}

Articolul sintetizează trei componente canonice ale proiectului: specificația
normativă, corpusul de bază și formalizarea gramaticală explicită. Specificația
oferă regulile active, corpusul furnizează exemplele validate și loturile
minimale de control, iar secțiunile analitice ale documentului motivează de ce
această propunere de normalizare este necesară și cum încearcă ea să reducă
costul flexiunii, al excepțiilor și al deducției contextuale fragile. În forma
publicată aici, aceste componente sunt tratate drept surse primare explicite ale
articolului, împreună cu repository-ul public al proiectului și cu versiunea în
engleză folosită ca punct comparativ de prezentare externă
\cite{rs-spec,rs-corpus,rs-eval,rs-prez,rs-readme,rs-repo,minenglish}.

\section{Arhitectura nucleului}

\subsection{Ordinea propoziției și tratamentul timpului}

Schema preferată a propoziției simple este \rs{subiect + timp + predicat}
cu continuări scurte. Fără marker temporal explicit, schema se reduce la
\rs{subiect + predicat}. Propoziția fără subiect exprimat rămâne permisă în
nucleul obligatoriu numai pentru impersonale reale, de tipul \rs{ploua.}

Timpul se scrie înaintea verbului. În uzul curent sunt preferate formele parțiale
sau relative, precum \rs{9}, \rs{20:30}, \rs{-1D} sau \rs{3H}. Forma completă
calendaristică rămâne rezervată stratului tehnic. Această decizie mută sarcina
temporală din flexiune într-un marker explicit și comparabil.

\begin{table}[h!]
\centering
\RStablelayout
\begin{tabularx}{\linewidth}{@{}lY@{}}
\rowcolor{RSGray}
\textbf{Funcție} & \textbf{Exemplu validat} \\
Timp punctual parțial & \rs{eu 9 mergea.} \\
Timp punctual cu minute & \rs{tu 20:30 vedea @film?} \\
Timp relativ trecut & \rs{eu -1D mergea la film.} \\
Timp relativ viitor & \rs{ei 3H venea.} \\
Impersonală fără subiect & \rs{ploua.} \\
\bottomrule
\end{tabularx}
\caption{Exemple minimale pentru ordinea preferată și markerii temporali}
\end{table}

\subsection{Forma verbală canonică}

Nucleul folosește o formă verbală stabilă, independentă de persoană și număr.
Pentru majoritatea verbelor active, forma canonică este imperfectul de persoana
a treia singular, tratat ca o convenție lexicală fixă. Copula reprezintă o
excepție controlată: \rs{a fi -> e}. Motivația este pragmatică: \rs{e} este mai
natural pentru cititorul român și mai ușor de procesat decât o copulă zero sau
decât generalizarea formei \rs{era}.

Această alegere nu rezolvă complet problema reziduului temporal al formei
canonice. Tocmai de aceea, proiectul o tratează ca ipoteză de lucru controlată,
nu ca adevăr teoretic închis. Timpul principal trebuie recuperat din markerul
temporal sau din context, nu din flexiunea istorică a verbului.

\subsection{Lexicul verbal și stratificarea sa}

Pentru a evita extinderea necontrolată a sistemului, verbele sunt împărțite în
trei straturi: \textit{nucleu obligatoriu}, \textit{permis, dar nerecomandat} și
\textit{lexic extins}. Această separare reduce presiunea de a promova în nucleu
orice verb lizibil, dar rar.

\begin{table}[h!]
\centering
\RStablelayout
\begin{tabularx}{\linewidth}{@{}lY@{}}
\rowcolor{RSGray}
\textbf{Strat} & \textbf{Inventar reprezentativ} \\
Nucleu obligatoriu & \rs{mergea}, \rs{vedea}, \rs{bea}, \rs{venea}, \rs{citea}, \rs{pleca}, \rs{rămânea}, \rs{intra}, \rs{vorbea}, \rs{ploua}, \rs{e} \\
Permis, dar nerecomandat & \rs{dădea}, \rs{spunea}, \rs{ajungea}, \rs{întreba}, \rs{răspundea}, \rs{cerea} \\
Lexic extins & \rs{făcea}, \rs{credea}, \rs{avea}, \rs{dormea}, \rs{stătea}, \rs{punea}, \rs{privea}, \rs{striga} \\
\bottomrule
\end{tabularx}
\caption{Stratificarea curentă a lexicului verbal}
\end{table}

\subsection{Nucleul nominal: definitudine, posesie și cardinalitate}

Forma de bază a substantivului este nearticulată. Definitudinea se marchează
numai când este necesară, prin operatorul invariabil \rs{@}. Posesia se scrie
analitic, cu operatorul \rs{\#}, în schema \rs{posesor\#obiect}. Numărul este
preferențial explicit prin prefix numeric, iar aproximarea nominală prin
\rs{\textasciitilde{}}.

\begin{table}[h!]
\centering
\RStablelayout
\begin{tabularx}{\linewidth}{@{}lY@{}}
\rowcolor{RSGray}
\textbf{Fenomen} & \textbf{Exemplu validat} \\
Nedefinit / generic & \rs{film} \\
Definit & \rs{@film} \\
Cardinalitate exactă & \rs{3cal} \\
Aproximare & \rs{\textasciitilde10pasăre} \\
Posesie simplă & \rs{Ana\#@carte e aici.} \\
\bottomrule
\end{tabularx}
\caption{Marcajele nominale centrale ale nucleului}
\end{table}

\subsection{Pronume cuantificate}

Una dintre cele mai distinctive decizii ale versiunii curente este introducerea
pronumelor cuantificate în nucleul obligatoriu. Forme precum \rs{eu5}, \rs{tu3},
\rs{ei10} sau \rs{eu*} sunt acceptate atunci când cardinalitatea
participanților este parte centrală a mesajului. În această configurație,
\rs{*} marchează plural neexact.

Aceste forme sunt admise în două poziții: ca subiect și ca posesor compact.
Astfel, \rs{eu3\#carte} înseamnă \textit{cartea noastră, a celor trei}, iar
\rs{eu3\#5carte} comprimă în aceeași unitate posesia de grup și cardinalitatea
obiectului posedat. Regula rămâne strictă: asemenea clustere sunt tolerate doar
când comprimă o relație nominală centrală fără pierdere de informație.

\subsection{Economia operatorială}

Nucleul actual încearcă să rămână sub un prag de densitate vizuală. În loturile
de test se preferă, de regulă, cel mult două marcaje vizibile, fără a număra
markerul temporal și semnul final. Un cluster cuantificat, de tipul \rs{eu5} sau
\rs{ei*}, este tratat ca un singur marcaj vizibil.

Operatorii centrali ai nucleului sunt \rs{@}, \rs{\#}, \rs{\textasciitilde{}}, markerii
temporali și semnele finale \rs{?} / \rs{!}, plus formele intensificate
\rs{??}, \rs{!!}, \rs{<<} și \rs{>>} atunci când sursa cere explicit un efect
mai puternic. În schimb, coordonarea și alternativa sunt lexicalizate
preferențial prin \rs{și}, \rs{sau} și \rs{ori}; simbolurile \rs{\&}, \rs{\textbar{}}
și \rs{\textasciicircum{}} rămân doar abrevieri opționale ale stratului extins.

\section{Sistemul operatorial}

\subsection{De ce operatorii sunt centrali în proiect}

RomânăSimplă nu poate fi înțeleasă corect dacă operatorii sunt tratați doar ca
detalii grafice. Ei sunt chiar mecanismul prin care proiectul mută informație
din flexiune în marcaje explicite, scurte și repetabile. În româna standard,
definitudinea, posesia, numărul, gradul și o parte din orientarea enunțului sunt
împrăștiate între articol, sufixe, acord și intonație. În RomânăSimplă, o parte
din această informație este concentrată în operatori puțini, dar stabili.

Avantajul acestei soluții este controlul local al sensului. Cititorul nu mai
trebuie să deducă din mai multe schimbări simultane dacă un substantiv este
definit, dacă o relație este posesivă sau dacă o propoziție este interogativă.
Costul este vizibilitatea operatorilor. Din acest motiv, proiectul nu urmărește
maximizarea numărului de operatori, ci disciplinarea lor prin regula
\textit{operator minim necesar}.

\subsection{Inventarul operatorilor activi}

În forma actuală, operatorii nu au toți același statut. Unii aparțin nucleului
obligatoriu, alții rămân doar în stratul extins, iar câțiva sunt admiși numai ca
intensificări controlate.

\begin{table}[h!]
\centering
\RStablelayout
\begin{tabularx}{\linewidth}{@{}l l Y@{}}
\rowcolor{RSGray}
\textbf{Operator} & \textbf{Statut} & \textbf{Funcție principală} \\
\rs{@} & nucleu & marchează definitudinea numai când trebuie exprimată explicit \\
\rs{\#} & nucleu & marchează posesia sau apartenența în schema \rs{posesor\#obiect} \\
\rs{\textasciitilde{}} & nucleu & marchează aproximarea, mai ales în cardinalitate nominală \\
\rs{<}, \rs{>} & nucleu & marchează grad explicit sub nivelul de bază sau peste nivelul de bază \\
\rs{<<}, \rs{>>} & nucleu controlat & intensifică gradul fără a deschide operatori noi \\
\rs{?}, \rs{!} & nucleu & marchează tipul final al enunțului \\
\rs{??}, \rs{!!} & nucleu controlat & intensifică interogația sau exclamarea când sursa cere contrast \\
markerii temporali & nucleu & mută timpul înaintea verbului, în formă parțială sau relativă \\
\rs{și}, \rs{sau}, \rs{ori} & extins preferat & coordonare și alternativă explicită în formă lexicală \\
\rs{\&}, \rs{\textbar{}}, \rs{\textasciicircum{}} & extins abreviat & prescurtări simbolice pentru stratul extins \\
\rs{\{} \rs{\}} & extins & delimitare pentru conținut raportat \\
\bottomrule
\end{tabularx}
\caption{Inventarul operatorial și statutul său curent}
\end{table}

Tabelul arată structura reală a proiectului: nucleul nu este construit pe o
abundență de simboluri, ci pe un grup mic de operatori care rezolvă problemele
cele mai costisitoare ale morfologiei standard. Extinsul păstrează explicit ceea
ce este util, dar încă neesențial pentru nucleul obligatoriu.

\subsection{Operatorii nominali: \rs{@}, \rs{\#}, număr și \rs{*}}

În zona nominală, operatorii sunt coloana vertebrală a simplificării.
\rs{@} separă forma lexicală a substantivului de definitudine. \rs{\#} scoate
posesia din genitiv și o face liniară. Prefixul numeric face cardinalitatea
imediat vizibilă, iar \rs{*} oferă o soluție economică pentru pluralul
nenumărat sau nedeterminat.

\begin{longtable}{@{}p{3.2cm}p{4.7cm}p{4.2cm}@{}}
\rowcolor{RSGray}
\textbf{Funcție} & \textbf{Formă RomânăSimplă} & \textbf{Lectură orientativă} \\
definitudine minimă & \rs{@film} & filmul \\
cardinalitate exactă & \rs{2cal} & doi cai \\
plural nedefinit & \rs{*pasăre} & multe păsări / păsări în general \\
definit plus număr & \rs{@2cal} & cei doi cai \\
posesie simplă & \rs{Ana\#@casă} & casa Anei \\
posesie cu obiect numerotat & \rs{eu\#2cal} & cei doi cai ai mei \\
posesie de grup compactă & \rs{eu3\#carte} & cartea noastră, a celor trei \\
posesie de grup plus cardinalitate & \rs{eu3\#5carte} & cele cinci cărți ale noastre, ale celor trei \\
posesie cu plural neexact personal & \rs{ei*\#@sală} & sala lor, a grupului nespecificat \\
\bottomrule
\caption{Operatorii nominali în uzul canonic și în compactarea extinsă}
\end{longtable}

Aceste exemple arată și una dintre tensiunile proiectului. Operatorii nominali
sunt foarte eficienți când exprimă exact o relație centrală. Devine însă riscant
dacă se acumulează prea multe informații periferice în aceeași unitate. De aici
vine diferența dintre o formă defensabilă ca \rs{eu3\#5carte} și o formă prea
încărcată ca \rs{eu5 vedea @2cal >clar.}, care amestecă într-o propoziție scurtă
mai multe tipuri de marcaj decât cere fenomenul dominant.

\subsection{Operatorii de grad și de enunț}

Operatorii \rs{<}, \rs{>}, \rs{\textasciitilde{}}, \rs{?} și \rs{!}
controlează o altă zonă sensibilă: gradul, aproximația și tipul propoziției.
Proiectul refuză să trateze gradul pozitiv simplu ca pe un caz marcat. Astfel,
\rs{@casă e mare.} este forma de bază, iar \rs{@casă e >mare.} apare doar când
surplusul de grad trebuie spus explicit.

Formele duble \rs{<<}, \rs{>>}, \rs{??} și \rs{!!} sunt importante teoretic,
pentru că ele arată cum proiectul încearcă să păstreze o familie operatorială
mică. Dublarea nu înseamnă operator nou, ci doar intensificarea unuia deja
admis. Această regulă menține sistemul compact și previne proliferarea de semne
cu funcții aproape identice.

\begin{table}[h!]
\centering
\RStablelayout
\begin{tabularx}{\linewidth}{@{}lY@{}}
\rowcolor{RSGray}
\textbf{Exemplu} & \textbf{Observație} \\
\rs{@casă e mare.} & adjectiv pozitiv simplu, fără grad redundant \\
\rs{@casă e >mare.} & grad explicit peste nivelul de bază \\
\rs{@casă e >>mare.} & intensificare reală a gradului, nu categorie nouă \\
\rs{eu vedea \textasciitilde10pasăre.} & aproximație numerică locală \\
\rs{tu 3H bea vin?} & interogație simplă \\
\rs{tu 3H bea vin??} & interogație intensificată \\
\rs{Ana vorbea!} & exclamare simplă \\
\rs{Ana vorbea!!} & exclamare intensificată \\
\bottomrule
\end{tabularx}
\caption{Operatorii de grad și de enunț în exemple scurte}
\end{table}

\subsection{Coordonare, alternativă și simbolizare controlată}

O decizie importantă a proiectului este că relațiile de coordonare și de alegere
nu mai sunt împinse în nucleul obligatoriu prin simboluri. Formele principale sunt
lexicale: \rs{și} pentru conjuncție, \rs{sau} pentru OR incluziv și \rs{ori}
pentru alternativă exclusivă. Simbolurile \rs{\&}, \rs{\textbar{}} și
\rs{\textasciicircum{}} rămân doar abrevieri opționale în stratul extins.

Această alegere este importantă pentru publicare, fiindcă ea arată că proiectul
nu urmărește tehnicizarea limbii cu orice preț. Acolo unde româna are deja forme
lexicale scurte și stabile, acestea sunt preferate. Operatorii simbolici sunt
rezervați pentru locurile unde scad clar costul structural.

\begin{longtable}{@{}p{3.6cm}p{4.2cm}p{4.4cm}@{}}
\rowcolor{RSGray}
\textbf{Relație} & \textbf{Formă preferată} & \textbf{Formă abreviată, doar extinsă} \\
Coordonare simplă & \rs{Ana citea și tu scria.} & \rs{Ana citea \& tu scria.} \\
Alternativă incluzivă & \rs{tu bea apă sau tu bea lapte.} & \rs{tu bea apă \textbar{} tu bea lapte.} \\
Alternativă exclusivă & \rs{tu bea apă ori tu bea lapte.} & \rs{tu bea apă \textasciicircum{} tu bea lapte.} \\
\bottomrule
\caption{Formele lexicale și simbolice pentru coordonare și alternativă}
\end{longtable}

\subsection{Buget operatorial și praguri de acceptare}

Bugetul operatorial este una dintre cele mai originale, dar și cele mai severe
reguli ale proiectului. El spune, în practică, că nu este suficient ca o formă
să fie corectă local; ea trebuie să rămână și lizibilă global. În nucleul de
test, propoziția preferă în mod normal cel mult două marcaje vizibile, fără a
număra markerul temporal și semnul final. Un cluster cuantificat este numărat ca
un singur marcaj.

\begin{table}[h!]
\centering
\RStablelayout
\begin{tabularx}{\linewidth}{@{}YlY@{}}
\rowcolor{RSGray}
\textbf{Exemplu} & \textbf{Verdict} & \textbf{Motiv} \\
\rs{eu5 3H venea.} & acceptat & cluster cuantificat plus timp; încă lizibil și localizat \\
\rs{eu5 vedea @2cal.} & acceptat & cardinalitatea subiectului și un singur marcaj nominal central \\
\rs{eu3\#5carte.} & acceptat condiționat & dens, dar comprimă exact relația posesivă centrală \\
\rs{eu5 vedea @2cal >clar.} & respins din nucleul de test & prea multe marcaje simultane pentru un singur enunț scurt \\
\bottomrule
\end{tabularx}
\caption{Exemple canonice pentru controlul bugetului operatorial}
\end{table}

Această grilă este utilă și editorial: ea oferă un criteriu explicit prin care
pot fi separate propozițiile bune pentru nucleul curent de propozițiile doar
ingenioase. Într-un context de publicare, regula aceasta este importantă pentru
că dă proiectului o disciplină verificabilă, nu doar un gust stilistic.

\section{Exemple și studii de caz}

\subsection{Secțiune extinsă de exemple comparative side by side}

Pentru publicare, exemplele nu trebuie lăsate la nivel de inventar scurt sau de
probă izolată.
Secțiunea de mai jos adună într-un singur loc o serie extinsă de perechi
comparative \textit{side by side}, construite astfel încât să arate nu doar
fenomene izolate, ci și combinații controlate de timp, posesie, cardinalitate,
copulă, relații prepoziționale, grad, aproximație, interogație, descompunere și
strat extins. Scopul nu este spectaculozitatea expresivă, ci densitatea
comparativă: cititorul trebuie să poată vedea repede cum se mută aceeași
informație din româna standard într-o schemă mai regulată.

\subsubsection{I. Nucleu temporal și propozițional}

\RSDenseTablelayout
\begin{longtable}{@{}p{0.5cm}p{5.45cm}p{5.45cm}@{}}
\rowcolor{RSGray}
\textbf{Nr.} & \textbf{Română standard} & \textbf{RomânăSimplă} \\
1 & Eu merg la ora 9 spre magazin. & \rs{eu 9 mergea la magazin.} \\
2 & Tu vezi filmul la 20:30? & \rs{tu 20:30 vedea @film?} \\
3 & Ei vin peste trei ore din sat. & \rs{ei 3H venea din sat.} \\
4 & Eu am văzut filmul ieri. & \rs{eu -1D vedea @film.} \\
5 & Ana citește cartea la ora 8. & \rs{Ana 8 citea @carte.} \\
6 & Tu vorbești clar cu Ana. & \rs{tu vorbea clar cu Ana.} \\
7 & Ei beau apă târziu. & \rs{ei bea apă târziu.} \\
8 & Tu bei vin peste trei ore? & \rs{tu 3H bea vin?} \\
9 & Tu bei vin peste trei ore, chiar? & \rs{tu 3H bea vin??} \\
10 & Ana vorbește foarte tare! & \rs{Ana vorbea!!} \\
11 & Plouă târziu. Noi rămânem acasă. & \rs{ploua târziu. eu rămânea acasă.} \\
12 & Ei vin târziu. Sala este închisă. & \rs{ei venea târziu. sală e închisă.} \\
13 & Ana intră în sală la ora 7. & \rs{Ana 7 intra în sală.} \\
14 & Ei pleacă peste două zile din sat. & \rs{ei 2D pleca din sat.} \\
15 & Eu citesc la ora 6. & \rs{eu 6 citea.} \\
16 & Tu intri peste treizeci de minute. & \rs{tu 30m intra.} \\
17 & Ei vin mâine la magazin. & \rs{ei 1D venea la magazin.} \\
18 & Ana vorbește repede după prânz. & \rs{Ana vorbea repede după.} \\
\bottomrule
\caption{Exemple comparative pentru nucleul temporal și propozițional}
\end{longtable}

\subsubsection{II. Relații nominale, cardinalitate și operatori}

\RSDenseTablelayout
\begin{longtable}{@{}p{0.5cm}p{5.45cm}p{5.45cm}@{}}
\rowcolor{RSGray}
\textbf{Nr.} & \textbf{Română standard} & \textbf{RomânăSimplă} \\
19 & Cartea Anei este aici. & \rs{Ana\#@carte e aici.} \\
20 & Casa Anei este mare. & \rs{Ana\#@casă e mare.} \\
21 & Trei cai intră în sală. & \rs{3cal intra în sală.} \\
22 & Eu văd aproximativ zece păsări. & \rs{eu vedea ~10pasăre.} \\
23 & Noi trei venim peste două zile. & \rs{eu3 2D venea.} \\
24 & Noi, mai mulți, rămânem acasă. & \rs{eu* rămânea acasă.} \\
25 & Cei doi cai ai mei intră. & \rs{eu\#@2cal intra.} \\
26 & Cele cinci cărți ale noastre sunt aici. & \rs{eu3\#5carte e aici.} \\
27 & Sala lor este închisă. & \rs{ei*\#@sală e închisă.} \\
28 & Fructul este copt. & \rs{@fruct e copt.} \\
29 & Casa este foarte mare. & \rs{@casă e >>mare.} \\
30 & Drumul este puțin închis. & \rs{@drum e <închis.} \\
31 & Maria vede cei doi cai. & \rs{Maria vedea @2cal.} \\
32 & Ion citește cartea mea. & \rs{eu\#@carte citea Ion.} \\
33 & Ușa casei este aici. & \rs{@casă\#@ușă e aici.} \\
34 & Noi trei vedem filmul. & \rs{eu3 vedea @film.} \\
35 & Ei, mai mulți, văd păsările. & \rs{ei* vedea @*pasăre.} \\
36 & Ana are cinci cărți. & \rs{Ana\#5carte.} \\
\bottomrule
\caption{Exemple comparative pentru relații nominale și operatori}
\end{longtable}

\subsubsection{III. Strat extins și cazuri de frontieră}

\RSDenseTablelayout
\begin{longtable}{@{}p{0.5cm}p{5.45cm}p{5.45cm}@{}}
\rowcolor{RSGray}
\textbf{Nr.} & \textbf{Română standard} & \textbf{RomânăSimplă} \\
37 & Ana citește și tu vorbești. & \rs{Ana citea și tu vorbea.} \\
38 & Tu bei apă sau lapte. & \rs{tu bea apă sau tu bea lapte.} \\
39 & Tu bei apă ori lapte, nu ambele. & \rs{tu bea apă ori tu bea lapte.} \\
40 & El spune: tu bei vin? & \rs{el spunea \{tu bea vin?\}} \\
41 & Ea crede că ei ajung la sală peste o zi. & \rs{ea credea \{ei 1D ajungea la sală.\}} \\
42 & Eu merg cu Ana la magazin peste trei ore. & \rs{eu 3H mergea cu Ana la magazin.} \\
43 & Ei dau apă la cal. & \rs{ei dădea apă la cal.} \\
44 & Ana spune că Maria vine. & \rs{Ana spunea \{Maria venea.\}} \\
45 & Tu întrebi și eu răspund după. & \rs{tu întreba și eu răspundea după.} \\
46 & Ei ajung la sală peste două ore. & \rs{ei 2H ajungea la sală.} \\
47 & Eu cer apă pentru Ana. & \rs{eu cerea apă pentru Ana.} \\
48 & El crede că filmul este bun. & \rs{el credea \{@film e mare.\}} \\
49 & Ana face pâine acasă. & \rs{Ana făcea pâine acasă.} \\
50 & Ion dormește în casă. & \rs{Ion dormea în casă.} \\
\bottomrule
\caption{Exemple comparative pentru stratul extins și pentru cazuri de frontieră}
\end{longtable}

\subsection{Micro-comparații de compactare}

Nu toate exemplele devin spectaculos mai scurte. În cazul românei, câștigul nu
vine întotdeauna din reducerea brută a numărului de caractere, ci din mutarea
informației într-o schemă mai regulată. Exemplele de mai jos sunt utile tocmai
pentru că arată unde proiectul câștigă mai ales în analizabilitate și unde
câștigă și în compresie locală.

\begin{table}[h!]
\centering
\RStablelayout
\begin{tabularx}{\linewidth}{@{}Y Y Y@{}}
\rowcolor{RSGray}
\textbf{Română standard} & \textbf{RomânăSimplă} & \textbf{Observație} \\
Eu merg la ora 9. & \rs{eu 9 mergea.} & timp mai explicit și formă verbală stabilă \\
Cartea Anei este aici. & \rs{Ana\#@carte e aici.} & posesie și definitudine liniarizate \\
Tu vezi filmul? & \rs{tu vedea @film?} & interogație și definitudine marcate local \\
Noi trei avem cinci cărți. & \rs{eu3\#5carte.} & compresie puternică, dar cu densitate mare \\
\bottomrule
\end{tabularx}
\caption{Micro-comparații între textul standard și forma simplificată}
\end{table}

\subsection{Cazuri-limită și praguri editoriale}

Pentru o versiune publicabilă, este important să nu fie ascunse exemplele în care
proiectul devine tensionat. Cazurile-limită nu sunt accidente; ele arată unde se
află marginea de siguranță a sistemului.

Primul tip de caz-limită este propoziția prea densă operatorial. Al doilea este
propoziția în care compactarea împinge cititorul să reconstruiască prea mult din
context. Al treilea este propoziția în care o relație ar merita descompusă, dar
este păstrată într-o singură unitate din dorință de compresie maximă. În toate
aceste situații, proiectul actual preferă prudența și coboară forma din nucleu
sau o rescrie.

\section{Corpus și validare}

\subsection{Perechi de control local}

Corpusul de bază este organizat în loturi scurte, fiecare construit pentru a
izola un fenomen dominant. Perechile comparative \textit{RO/SR} evită amestecul
de reguli și permit verificarea locală a deciziilor de transcriere.

\begin{longtable}{@{}p{0.75cm}p{5.35cm}p{6.1cm}@{}}
\rowcolor{RSGray}
\textbf{Nr.} & \textbf{Română standard} & \textbf{RomânăSimplă} \\
1 & Eu merg la ora 9. & \rs{eu 9 mergea.} \\
2 & Tu vezi filmul. & \rs{tu vedea @film.} \\
3 & Ana intră în sală. & \rs{Ana intra în sală.} \\
4 & Eu văd aproximativ zece păsări. & \rs{eu vedea ~10pasăre.} \\
5 & Cartea Anei este aici. & \rs{Ana\#@carte e aici.} \\
6 & Trei cai intră. & \rs{3cal intra.} \\
7 & Casa este mare. & \rs{@casă e mare.} \\
8 & Plouă. Noi rămânem acasă. & \rs{ploua. eu rămânea acasă.} \\
9 & Tu vorbești clar. & \rs{tu vorbea clar.} \\
10 & Ei pleacă. & \rs{ei pleca.} \\
\bottomrule
\caption{Eșantion de control din corpusul canonic de bază}
\end{longtable}

Această structură urmărește două obiective. Mai întâi, fiecare propoziție trebuie
să rămână lizibilă la prima lectură. Apoi, fiecare regulă trebuie să poată fi
testată repetabil fără a depinde de un exemplu izolat sau de o justificare
intuitivă post hoc.

\subsection{Loturi specifice de control}

Pe lângă perechile comparative, proiectul publică loturi dedicate de control pentru timp,
invarianță verbală, structură, număr, definitudine și plural uman. De pildă,
lotul de plural uman fixează explicit atât subiectele cuantificate, cât și
posesia compactă de grup: \rs{eu2 venea.}, \rs{eu* rămânea acasă.},
\rs{eu3\#carte.}, \rs{eu3\#5carte.}, \rs{ei*\#@sală.}

În paralel, controlul rapid verifică patru lucruri: ordinea preferată a
propoziției, forma verbală stabilă, forma nemarcată păstrată unde este suficientă
și aplicarea regulii \textit{operator minim necesar}. Din punct de vedere
metodologic, aceasta este una dintre cele mai puternice părți ale proiectului:
corpusul nu conține doar exemple acceptate, ci și exemple ambigue, dense sau
respinse.

\section{Motivație și justificare teoretică}

\subsection{Ce justifică propunerea}

Propunerea de normalizare devine defensabilă tocmai pentru că nucleul actual
este mic, coerent și documentat prin loturi separate. Patru puncte sunt deja
solide ca justificare de sistem.

\begin{enumerate}[leftmargin=*]
\item Mutarea timpului înaintea verbului scade povara paradigmelor verbale și
permite reprezentări compacte și predictibile.
\item Definitudinea prin \rs{@} și posesia prin \rs{\#} sunt scurte, stabile și
ușor de segmentat.
\item Descompunerea frazelor grele reduce nevoia de imbricare și crește
lizibilitatea comparativă.
\item Stratificarea lexicului verbal împiedică diluarea nucleului prin expansiune
lexicală oportunistă.
\end{enumerate}

Punctele cele mai solide apar în zona de validare și reproductibilitate. Asta
înseamnă că proiectul nu mai depinde exclusiv de impresii stilistice, ci oferă
un aparat compact de control pe care un evaluator extern îl poate reface.

\subsection{Limite teoretice majore}

Două probleme rămân centrale. Prima este forma verbală canonică. Deși sistemul
declară că timpul vine din marker și context, forma de imperfect păstrează un
reziduu istoric care nu poate fi eliminat prin decret. Situația este gestionată
onest prin loturi de invarianță verbală și prin comparație explicită cu alte
forme candidate, dar nu este complet rezolvată.

A doua problemă este presiunea viitoare de promovare lexicală. Orice limbă
simplificată riscă să își piardă coerența exact atunci când începe să accepte,
treptat, prea multe excepții sau prea multe verbe motivate local. De aceea,
versiunea actuală insistă pe un nucleu verbal înghețat și pe o zonă \textit{permis,
dar nerecomandat} clar separată.

\subsection{Densitate operatorială și compactare}

Introducerea pronumelor cuantificate și acceptarea unor forme dense, precum
\rs{eu3\#5carte.}, aduc un câștig real de compactare. În același timp, ele arată
limita exactă la care compactarea poate începe să concureze cu lizibilitatea.
Soluția actuală este prudentă: clusterele dense rămân acceptabile doar când ele
comprimă o relație centrală, nu când acumulează marcaje de dragul simetriei.

În acest sens, proiectul evită o formulă maximalistă. Nu orice propoziție scurtă
este și bună; bună este numai propoziția scurtă care rămâne imediat parsabilă.

\subsection{De ce secțiunea despre operatori este decisivă pentru publicare}

Pentru o lectură academică externă, cel mai vulnerabil punct al proiectului nu
este neapărat forma verbală, ci riscul ca operatorii să pară o adunare
eclectică de convenții. Tocmai de aceea, o prezentare publicabilă trebuie să
arate coerent familia operatorială: ce intră în nucleu, ce rămâne extins, cum se
numără densitatea și de ce formele duble nu reprezintă operatori noi.

Dacă acest cadru lipsește, proiectul poate fi citit nedrept ca o simplă
romanizare tehnicizată. Dacă însă operatorii sunt prezentați ca un sistem cu
inventar limitat, cu praguri de acceptare și cu exemple validate, atunci nucleul
devine mult mai ușor de apărat în fața unui evaluator extern.

\section{Discuție}

Forma actuală a RomânăSimplă este interesantă tocmai pentru că nu încearcă să
fie o limbă complet nouă, ci o varietate controlată a românei, construită cu
instrumente relativ puține. Sistemul reduce flexiunea, dar nu renunță la
ortografia standard; introduce operatori, dar încearcă să îi țină puțini;
admite compactarea, dar o limitează prin reguli explicite de densitate.

Această combinație îl face mai defensabil academic decât multe proiecte de
reformă totală. Totuși, aceeași alegere îl face și mai vulnerabil la obiecția
„cât de departe poate merge fără să își piardă naturalețea?”. Răspunsul actual
este unul deliberat modest: proiectul nu forțează încă un răspuns global, ci își
consolidează nucleul prin loturi de control și prin refuzul de a promova prea
repede reguli sau forme noi.

\section{Concluzie}

RomânăSimplă are, în starea actuală, un nucleu suficient de coerent pentru a
susține o prezentare academică serioasă. Cele mai solide decizii sunt tratamentul
timpului, economia nominală prin \rs{@} și \rs{\#}, separarea strictă a
straturilor lexicale și metodologia de validare pe loturi mici. Cele mai mari
rezerve rămân forma verbală canonică și disciplina necesară pentru a împiedica
extinderea necontrolată a sistemului.

Contribuția principală a proiectului nu este că oferă deja o gramatică totală,
ci că propune un nucleu minim de normalizare, cu criterii explicite de acceptare,
cu exemple verificabile și cu o formalizare suficient de clară pentru parsare,
revizie și extindere controlată. În acest sens, RomânăSimplă este mai aproape de
un protocol lingvistic de lucru și de o propunere de standardizare locală decât
de o reformă complet închisă.

\appendix

\section{Inventar scurt de forme reprezentative}

\begin{table}[h!]
\centering
\RStablelayout
\begin{tabularx}{\linewidth}{@{}lY@{}}
\rowcolor{RSGray}
\textbf{Categorie} & \textbf{Exemple scurte} \\
Timp & \rs{eu 9 mergea.}, \rs{ei -1D venea.}, \rs{Ana 3H citea.} \\
Definitudine & \rs{film}, \rs{@film}, \rs{@2cal} \\
Posesie & \rs{Ana\#@carte}, \rs{eu3\#carte}, \rs{ei*\#@sală} \\
Cardinalitate & \rs{3cal}, \rs{\textasciitilde10pasăre}, \rs{eu5} \\
Structură simplă & \rs{@casă e mare.}, \rs{tu vorbea clar.}, \rs{ploua.} \\
\bottomrule
\end{tabularx}
\caption{Inventar scurt de exemple reprezentative}
\end{table}

\section{Rezumat al deciziilor normative}

\begin{enumerate}[leftmargin=*]
\item Timpul se scrie înaintea verbului.
\item Forma verbală este stabilă și nu depinde de persoană sau număr.
\item Copula preferată în nucleu este \rs{e}.
\item Definitudinea se marchează prin \rs{@} numai când este necesară.
\item Posesia se marchează prin \rs{\#}.
\item Coordonările preferă formele lexicale \rs{și}, \rs{sau}, \rs{ori}; simbolurile rămân extinse.
\item Pronumele cuantificate sunt admise în nucleu ca subiect și ca posesor compact.
\item Se aplică regula \textit{operator minim necesar}.
\end{enumerate}

\section{EBNF Core}

Anexa aceasta reproduce gramatica de nucleu într-o formă compactă, tocmai pentru
ca cititorul tehnic să poată verifica direct relația dintre argumentul teoretic,
corpusul validat și formalizarea publicată.

\RSEBNFLayout
\begin{verbatim}
core-text             = core-sentence, { ws, core-sentence };

core-sentence         = simple-sentence, end-mark;

simple-sentence       = explicit-subject-sentence
                      | nominal-sentence
                      | reduced-sentence;

nominal-sentence      = nominal-group;

explicit-subject-sentence	= subject, ws,
						[ time-marker, ws ],
						predicate;

reduced-sentence      = impersonal-predicate;

impersonal-predicate  = impersonal-verb, { ws, adverb };

predicate             = verbal-predicate
                      | copular-predicate;

verbal-predicate      = canonical-verb, { ws, continuation };

copular-predicate     = copula, ws, nonverbal-predicate;

nonverbal-predicate   = nominal-group
                      | prepositional-group
                      | adverb
                      | adjective-tail
                      | approximation-group;

continuation          = nominal-group
                      | demonstrative-atom
                      | prepositional-group
                      | adverb
                      | graded-adjective
                      | approximation-group;

subject               = pronoun
                      | quantified-pronoun
                      | proper-name
                      | nominal-group;

nominal-group         = [ possessor, "#" ], nominal-core, [ ws, adjective-tail ];
possessor             = pronoun | quantified-pronoun | proper-name | noun;
nominal-core          = [ "@" ], [ integer | "*" ], noun;
adjective-tail        = adjective | graded-adjective;
graded-adjective      = grade-operator, adjective;
approximation-group   = "~", integer, noun;

prepositional-group   = preposition, ws, nominal-group;
demonstrative-atom    = "asta";

time-marker           = clock-time | relative-time;
clock-time            = hour | hour, ":", minute;
relative-time         = [ "+" | "-" ], integer, time-unit;

grade-operator        = "<" | ">" | "<<" | ">>";
end-mark              = "." | "?" | "!" | "??" | "!!";

preposition           = "la" | "din" | "cu" | "pentru" | "pe" | "în";
time-unit             = "m" | "H" | "D";
canonical-verb        = lexical-verb;
copula                = "e";
impersonal-verb       = lexical-impersonal-verb;

pronoun               = "eu" | "tu" | "el" | "ea" | "ei";
quantified-pronoun    = pronoun, ( integer | "*" );

ws                    = " ", { " " };
integer               = digit, { digit };
digit                 = "0" | "1" | "2" | "3" | "4" | "5" | "6" | "7" | "8" | "9";
hour                  = digit, [ digit ];
minute                = digit, digit;
second                = digit, digit;
day                   = digit, digit;
month                 = digit, digit;
year                  = digit, digit, digit, digit;

proper-name           = lexical-proper-name;
noun                  = lexical-noun;
adjective             = lexical-adjective;
adverb                = lexical-adverb;
lexical-verb          = lexical-core-verb;
\end{verbatim}
\normalsize

\section{EBNF Extended și lexic minim}

\RSEBNFLayout
\begin{verbatim}
extended-sentence       = xor-sentence
                          | lexical-and-sentence
                          | lexical-or-sentence
                          | lexical-xor-sentence
                          | and-sentence
                          | or-sentence
                          | reported-sentence
                          | technical-time-sentence
                          | negated-verbal-sentence
                          | negated-copular-sentence
                          | negated-impersonal-sentence;

lexical-and-sentence    = simple-sentence, ws, "și", ws, simple-sentence, end-mark;
lexical-or-sentence     = simple-sentence, ws, "sau", ws, simple-sentence, end-mark;
lexical-xor-sentence    = simple-sentence, ws, "ori", ws, simple-sentence, end-mark;

and-sentence            = simple-sentence, ws, "&", ws, simple-sentence, end-mark;
or-sentence             = simple-sentence, ws, "|", ws, simple-sentence, end-mark;
xor-sentence            = simple-sentence, ws, "^", ws, simple-sentence, end-mark;

reported-sentence       = reporting-clause, ws, "{", embedded-sentence, "}";
reporting-clause        = [ subject, ws ], [ time-marker, ws ], reporting-verb;
embedded-sentence       = simple-sentence, end-mark;

technical-time-sentence = [ subject, ws ], technical-time, ws, predicate, end-mark;
technical-time          = year, "-", month, "-", day, ws,
                          hour, ":", minute, ":", second;

negated-verbal-sentence = subject, ws, [ time-marker, ws ], "!", canonical-verb,
                                                    { ws, continuation }, end-mark;
negated-copular-sentence	= subject, ws, [ time-marker, ws ],
						copula, ws, "!",
						( nominal-group
						| prepositional-group
						| adverb
						| adjective-tail
						| approximation-group ),
						end-mark;
negated-impersonal-sentence	= "!", impersonal-verb,
						{ ws, adverb }, end-mark;

reporting-verb          = "spunea" | "întreba" | "credea";

lexical-core-verb       = "mergea" | "vedea" | "bea" | "venea"
                          | "citea" | "pleca" | "rămânea"
                          | "intra" | "vorbea" | "ploua" | "e";

lexical-permitted-verb  = "dădea" | "spunea" | "ajungea"
                          | "întreba" | "răspundea" | "cerea";

lexical-extended-verb   = "făcea" | "credea" | "avea" | "dormea"
                          | "stătea" | "punea" | "privea" | "striga";

lexical-impersonal-verb = "ploua";
lexical-adverb          = "clar" | "repede" | "aici" | "acasă"
                          | "afară" | "târziu" | "după";
lexical-adjective       = "mare" | "copt" | "închis" | "gol" | "liniștit";
lexical-proper-name     = "Ana" | "Maria" | "Ion";
lexical-noun            = "film" | "carte" | "pasăre" | "casă" | "cal"
                          | "sală" | "magazin" | "vin" | "apă"
                          | "lapte" | "sat" | "ușă" | "pâine"
                          | "fruct" | "frig" | "drum" | "om";
\end{verbatim}
\normalsize

\section{Luceafărul: anexă ilustrativă}

Materialul de mai jos nu are statut normativ. El este păstrat în anexă tocmai
pentru a arăta cum se comportă sistemul atunci când este scos din zona de propoziții
scurte de control și este pus la lucru pe un text literar cu densitate imagistică și
ritm propriu. Această anexă este utilă editorial: ea arată ce poate face proiectul,
dar și unde începe să plătească prețul pentru compactare sau parafrază. Textul
poetic de referință folosit aici este \textit{Luceafărul} de Mihai Eminescu
\cite{eminescu-luceafarul}.

\small

\subsection*{Strofa 1}
\noindent
\begin{tabularx}{\linewidth}{@{}Y Y@{}}
\rowcolor{RSGray}
\textbf{Original} & \textbf{Adaptare SR} \\
\begin{minipage}[t]{\linewidth}\vspace{0pt}
A fost odată ca-n povești,\\
A fost ca niciodată,\\
Din rude mari împărătești,\\
O prea frumoasă fată.
\end{minipage}
&
\begin{minipage}[t]{\linewidth}\vspace{0pt}
\rs{a fost odată ca în *poveste.}\\
\rs{a fost ca niciodată.}\\
\rs{din *rudă >împărătesc,}\\
\rs{e o fată >frumoasă.}
\end{minipage}
\\
\end{tabularx}

\subsection*{Strofa 2}
\noindent
\begin{tabularx}{\linewidth}{@{}Y Y@{}}
\rowcolor{RSGray}
\textbf{Original} & \textbf{Adaptare SR} \\
\begin{minipage}[t]{\linewidth}\vspace{0pt}
Și era una la părinți\\
Și mândră-n toate cele,\\
Cum e Fecioara între sfinți\\
Și luna între stele.
\end{minipage}
&
\begin{minipage}[t]{\linewidth}\vspace{0pt}
\rs{și e una la *părinte.}\\
\rs{și mândră în toate cele,}\\
\rs{ea e ca Fecioara între *sfânt.}\\
\rs{și luna între *stea.}
\end{minipage}
\\
\end{tabularx}

\subsection*{Strofa 3}
\noindent
\begin{tabularx}{\linewidth}{@{}Y Y@{}}
\rowcolor{RSGray}
\textbf{Original} & \textbf{Adaptare SR} \\
\begin{minipage}[t]{\linewidth}\vspace{0pt}
Din umbra falnicelor bolți\\
Ea pasul și-l îndreaptă\\
Lângă fereastră, unde-n colț\\
Luceafărul așteaptă.
\end{minipage}
&
\begin{minipage}[t]{\linewidth}\vspace{0pt}
\rs{din umbra de sub *boltă înaltă,}\\
\rs{ea își ducea pasul.}\\
\rs{la fereastră, unde în colț,}\\
\rs{Luceafărul așteaptă.}
\end{minipage}
\\
\end{tabularx}

\subsection*{Strofa 4}
\noindent
\begin{tabularx}{\linewidth}{@{}Y Y@{}}
\rowcolor{RSGray}
\textbf{Original} & \textbf{Adaptare SR} \\
\begin{minipage}[t]{\linewidth}\vspace{0pt}
Privea în zare cum pe mări\\
Răsare și străluce,\\
Pe mișcătoarele cărări\\
Corăbii negre duce.
\end{minipage}
&
\begin{minipage}[t]{\linewidth}\vspace{0pt}
\rs{ea privea în zare, pe *mare,}\\
\rs{cum el răsare și străluce.}\\
\rs{pe *cărare mișcătoare,}\\
\rs{el ducea *corabie neagră.}
\end{minipage}
\\
\end{tabularx}

\subsection*{Strofa 5}
\noindent
\begin{tabularx}{\linewidth}{@{}Y Y@{}}
\rowcolor{RSGray}
\textbf{Original} & \textbf{Adaptare SR} \\
\begin{minipage}[t]{\linewidth}\vspace{0pt}
Îl vede azi, îl vede mâni,\\
Astfel dorința-i gata;\\
El iar, privind de săptămâni,\\
Îi cade dragă fata.
\end{minipage}
&
\begin{minipage}[t]{\linewidth}\vspace{0pt}
\rs{ea vedea pe el. ea 1D vedea pe el.}\\
\rs{astfel, ea\#dorință e gata.}\\
\rs{iar el o privea de *săptămână.}\\
\rs{fata îi cade dragă.}
\end{minipage}
\\
\end{tabularx}

\subsection*{Strofa 6}
\noindent
\begin{tabularx}{\linewidth}{@{}Y Y@{}}
\rowcolor{RSGray}
\textbf{Original} & \textbf{Adaptare SR} \\
\begin{minipage}[t]{\linewidth}\vspace{0pt}
Cum ea pe coate-și răzima\\
Visând ale ei tâmple,\\
De dorul lui și inima\\
Și sufletu-i se împle.
\end{minipage}
&
\begin{minipage}[t]{\linewidth}\vspace{0pt}
\rs{cum ea stătea pe *cot,}\\
\rs{visând, cu ea\#tâmplă,}\\
\rs{de el\#dor, ea\#inimă,}\\
\rs{și ea\#suflet se umple.}
\end{minipage}
\\
\end{tabularx}

\normalsize

\begin{center}
\small\textbf{Tabel.} Primele șase strofe din \textit{Luceafărul}, în format side by side, cu vers întreg pe linie.
\end{center}

\section{Observații asupra anexei literare}

Exemplul din \textit{Luceafărul} este util tocmai pentru că nu confirmă doar
punctele tari ale sistemului. El arată și dificultățile reale: poezia cere ritm,
ambiguitate controlată, densitate semantică și o negociere continuă între
literalitate și parafrază. Într-un asemenea context, RomânăSimplă poate produce
forme interesante, dar nu întotdeauna naturale sau economic superioare. Din
acest motiv, anexa literară trebuie citită ca probă exploratorie, nu ca dovadă
normativă pentru nucleul canonic.

\section{Disponibilitate și reproducibilitate}

Lucrarea este reproductibilă la nivel documentar. Regulile normative, corpusul
de bază, exemplele comparative și sursa \LaTeX\ sunt păstrate în repository-ul
public al proiectului \cite{rs-repo}. Pentru o evaluare externă sau pentru o
depunere pe platforme academice precum arXiv, acest lucru este esențial: cititorul
poate confrunta direct afirmațiile articolului cu documentele canonice, fără să
depindă de o reconstrucție informală a istoriei proiectului.

\begin{thebibliography}{9}

\bibitem{rs-spec}
Dan Micșa.
\textit{Specificația Limbii RomânăSimplă}.
Document canonic de lucru, martie 2026.

\bibitem{rs-corpus}
Dan Micșa.
\textit{Corpusul de Bază RomânăSimplă}.
Corpus canonic validat pentru nucleul curent, martie 2026.

\bibitem{rs-eval}
Dan Micșa.
\textit{Evaluare Academică}.
Document intern al proiectului RomânăSimplă, martie 2026.

\bibitem{rs-prez}
Dan Micșa.
\textit{Prezentare}.
Document intern de sinteză pentru proiectul RomânăSimplă, martie 2026.

\bibitem{rs-readme}
Dan Micșa.
\textit{README}.
Document de orientare al proiectului RomânăSimplă, martie 2026.

\bibitem{rs-repo}
Dan Micșa.
\textit{RomânăSimplă}.
Repository public GitHub: \url{https://github.com/dmicsa/Rom-n-Simpl-.git}.
Accesat în martie 2026.

\bibitem{minenglish}
Dan Micșa.
\textit{MinEnglish}.
Versiune în engleză în format PDF: \url{https://github.com/dmicsa/MinEnglish/blob/master/MinEnglish.pdf}.
Accesat în martie 2026.

\bibitem{eminescu-luceafarul}
Mihai Eminescu.
\textit{Luceafărul}.
Text poetic folosit în anexă exclusiv ca material ilustrativ comparativ.

\end{thebibliography}

\end{document}